\chapter{Dental Crown Procedure}
\index{Dental Crown Procedure}

\section*{Definition and Overview}
A dental crown is a cap-shaped restoration that fits over a damaged or weakened tooth to restore its shape, size, strength, and appearance. Crowns are used to protect a tooth after root canal treatment, to repair a tooth with a large filling or crack, and to improve the appearance of discoloured or misshapen teeth. The crown procedure is usually completed over two visits: preparing the tooth and taking impressions, then fitting the permanent crown. This chapter explains when crowns are needed, what the procedure involves, and how to care for a crowned tooth.

\section*{Epidemiology}
Crowns are among the most common restorative procedures in dentistry. Millions of crowns are placed worldwide each year, most commonly on molars and premolars with large fillings or after root canal treatment. The need for crowns increases with age as teeth accumulate damage, fillings, and endodontic treatment. Crowns are also used on dental implants. Materials have evolved from full metal to porcelain-fused-to-metal and all-ceramic options, with an increasing preference for tooth-coloured crowns for aesthetic reasons.

\section*{Aetiology and Pathophysiology}
\begin{itemize}
    \item \textbf{Large fillings:} a tooth with a very large filling is weakened and can fracture without full coverage
    \item \textbf{Root canal treatment:} endodontically treated teeth become brittle and need a crown for protection
    \item \textbf{Fractures and cracks:} trauma or wear can damage the tooth structure beyond simple filling repair
    \item \textbf{Decay:} extensive decay removes much of the natural tooth
    \item \textbf{Aesthetics:} crowns improve severely discoloured or misshapen teeth
    \item \textbf{Implant restoration:} crowns are attached to dental implants to replace missing teeth
\end{itemize}

\section*{Clinical Features}
\begin{itemize}
    \item A tooth that is heavily filled, cracked, or root-treated and at risk of further damage
    \item A tooth with a crown in place looks and functions like a natural tooth
    \item Problems to watch for: sensitivity, pain, a loose crown, or food trapping at the margin
    \item Gum irritation or recession around the crown
    \item Wear or fracture of the crown over time
\end{itemize}

\section*{Differential Diagnosis}
\begin{itemize}
    \item Crown versus filling: crowns cover the whole tooth and are used when a filling is insufficient
    \item Crown versus onlay or veneer: onlays cover part of the tooth, and veneers cover only the front surface
    \item Crown versus extraction and implant for severely damaged teeth
    \item Temporary versus permanent crown
\end{itemize}

\section*{When to Seek Medical Care}
A dentist will advise whether a crown is needed based on the condition of the tooth. See a dentist promptly if a crowned tooth becomes painful, sensitive, or loose, or if the crown chips or falls off.

\textbf{Contact your local emergency services for:}
\begin{itemize}
    \item Severe pain or swelling after the crown procedure, which may indicate infection
    \item Trauma that fractures the tooth or crown
    \item Uncontrolled bleeding from the treated area
\end{itemize}

\section*{Diagnosis and Investigations}
\begin{itemize}
    \item \textbf{Dental examination:} assessment of the tooth, existing restorations, and bite
    \item \textbf{X-rays:} check for decay, root health, and bone support
    \item \textbf{Assessment of the need:} whether a crown, onlay, or filling is most appropriate
    \item \textbf{Shade matching:} selection of the crown colour to match adjacent teeth
    \item \textbf{Impressions or digital scans:} records used by the laboratory to fabricate the crown
\end{itemize}

\section*{Management}
\begin{itemize}
    \item \textbf{First visit:} the tooth is numbed, shaped to receive the crown, and any decay is removed
    \item \textbf{Impressions:} an impression or digital scan is taken, and a temporary crown is fitted
    \item \textbf{Laboratory fabrication:} the permanent crown is made over 1--3 weeks
    \item \textbf{Second visit:} the temporary crown is removed and the permanent crown is tried, adjusted, and cemented
    \item \textbf{Same-day crowns:} CAD/CAM technology allows some crowns to be designed and milled in the surgery in one visit
    \item \textbf{Post-procedure care:} avoid hard foods until settled, and maintain good hygiene
    \item \textbf{Long-term review:} regular check-ups to monitor the crown and its margin
\end{itemize}

\section*{Complications}
\begin{itemize}
    \item Sensitivity to hot or cold after placement, which usually settles
    \item Decay or gum disease developing at the crown margin
    \item Loosening, chipping, or fracture of the crown
    \item Discomfort or sore gums around the crown
    \item Allergic reactions to metals in rare cases
    \item Need for root canal treatment if the tooth becomes symptomatic
    \item Replacement of the crown after many years of service
\end{itemize}

\section*{Prognosis and Outcomes}
With good oral hygiene and regular dental care, a well-made crown typically lasts 10--15 years or longer. Crowns restore the function and appearance of damaged teeth and protect them from further damage, and the vast majority of patients are satisfied with the result. The most common cause of crown failure is decay or gum disease at the margin, which is preventable with good home care and regular check-ups.

\section*{Prevention and Self-Care}
\begin{itemize}
    \item Brush twice daily and clean carefully around the crowned tooth
    \item Floss or use interdental brushes to keep the crown margin clean
    \item Attend regular check-ups so the crown and gums are monitored
    \item Avoid chewing very hard foods, ice, or other objects
    \item Wear a mouthguard if you grind your teeth or play contact sport
    \item Report sensitivity, pain, or a loose crown promptly
\end{itemize}

\section*{Sources and Further Reading}
The following reputable sources provide further information for healthcare students and patients seeking reliable, current advice about dental crowns.

\begin{itemize}
    \item Australian Dental Association. \textit{Dental crowns: information for patients.} https://www.ada.org.au/
    \item NHS. \textit{Dental crowns: procedure and aftercare.} https://www.nhs.uk/conditions/dental-crowns/
    \item Mayo Clinic. \textit{Dental crowns: what to expect.} https://www.mayoclinic.org/
    \item MedlinePlus. \textit{Dental crowns.} https://medlineplus.gov/
    \item Healthdirect Australia. \textit{Dental crowns and bridges.} https://www.healthdirect.gov.au/
    \item Oral Health Foundation. \textit{Crowns and bridges.} https://www.dentalhealth.org/
\end{itemize}
